\section{Widen before subtracting}
A stored reading lies in \([-2^{31},2^{31}-1]\). Subtracting two such readings can
give any difference between \(-(2^{32}-1)\) and \(2^{32}-1\). A signed 33-bit
integer is sufficient, while a signed word is not. For the largest upward jump,
\[
 (2^{31}-1)-(-2^{31})=4294967295.
\]
A word subtraction wraps that difference to \(-1\). Widening the result afterward
cannot distinguish that wrapped value from a real downward step of one.

Load each reading with signed \texttt{lw}, retaining its sign in the full register,
then subtract with \texttt{dsubu}. Despite its wraparound semantics, this doubleword
subtraction is exact throughout the proven input domain. The previous reading must
also retain its full signed value. An unsigned word load interprets stored negative
readings as large positive integers and changes the sequence before subtraction.

\begin{table}[htbp]
\centering
\caption{Differences whose widths are determined by the pair of readings.}
\tagpdfsetup{table/header-rows={1}}
\begin{tabular}{rrr}
\toprule
Previous & Current & Exact delta\\
\midrule
0 & \(-2147483648\) & \(-2147483648\)\\
0 & 2147483647 & 2147483647\\
\(-2147483648\) & 2147483647 & 4294967295\\
2147483647 & \(-2147483648\) & \(-4294967295\)\\
\bottomrule
\end{tabular}
\end{table}

\begin{MARSExercise}{Distinguish two encodings of a mistaken delta}
Trace an implementation that uses word subtraction for the upward endpoint jump.
Which delta does it pass to the signed mapping? Explain why a decoder that shares the
same mistake can appear to agree in a round-trip test while still violating this format.
\end{MARSExercise}
